% =====================================================================
% Prova formal: bound de energia H^1 para correntes Kac-Moody
% extraidas de uma MERA SU(2)_3-simetrica via residuos pull-through.
%
% Edivaldo Amaral - CS-MERA Series, Apendice tecnico.
% Versao: maio 2026.
%
% Dependencias: amsmath, amssymb, amsthm, mathtools.
% Ambientes esperados: theorem, lemma, proposition, definition,
%                     remark, proof (estilo "plain"/"definition").
% =====================================================================

\section{Bound de energia $H^1$ para correntes pull-through}
\label{sec:bound_H1}

Esta secção estabelece a estimativa de energia uniforme
\begin{equation}
  \|J_N^a(f)\,\psi\| \;\le\; C\,\|f\|_{C^1}\,\|\psi\|,
  \qquad C \text{ independente de } N,
  \label{eq:bound_H1_master}
\end{equation}
com constante explícita $C = 12$ para uma MERA binária e $C = 9$ para
uma MERA ternária, em ambos os casos com isometrias e desemaranhadores
no espaço de intertwinedores da categoria modular $\mathrm{SU}(2)_3$.
A prova procede por quatro lemas independentes que combinam-se em
forma multiplicativa.

\subsection{Notação e dados estruturais}

Fixamos as seguintes convenções, válidas em toda a secção.

\begin{description}
\item[Categoria.] $\mathcal{C} = \mathrm{SU}(2)_3$, com objectos simples
$j \in \{0, \tfrac{1}{2}, 1, \tfrac{3}{2}\}$ e dimensões quânticas
$d_0 = d_{3/2} = 1$, $d_{1/2} = d_1 = \varphi = (1+\sqrt{5})/2$.
Espaços vectoriais subjacentes:
\[
V_0 \cong \mathbb{C}, \qquad
V_{1/2} \cong \mathbb{C}^2, \qquad
V_1 \cong \mathbb{C}^3, \qquad
V_{3/2} \cong \mathbb{C}^4.
\]

\item[Geradores.] Para cada $j \in \{0, \tfrac{1}{2}, 1, \tfrac{3}{2}\}$,
denotamos por $t_j^a$ ($a=1,2,3$) a representação dos geradores de
$\mathfrak{su}(2)$ em $V_j$. Em $V_{1/2}$, $t^a = \sigma^a/2$, donde
$\|t^a_{1/2}\|_{\mathrm{op}} = 1/2$. Em geral, o autovalor máximo
de $t^z_j$ é $j$, donde
\begin{equation}
  \|t^a_j\|_{\mathrm{op}} \;=\; j, \qquad
  \max_{j \in \{0,1/2,1,3/2\}} \|t^a_j\|_{\mathrm{op}} \;=\; \tfrac{3}{2}.
  \label{eq:norma_ta}
\end{equation}
Definimos $\jmax := 3/2$.

\item[Tensores.] Os tensores locais $A$ da MERA pertencem a
\[
  \mathrm{Hom}_{\mathcal{C}}\!\left(\,\bigotimes_{i \in \mathrm{in}} V_{j_i},\;
  \bigotimes_{k \in \mathrm{out}} V_{j_k}\,\right).
\]
Numa MERA padrão com factor de coarse-graining $b \in \{2,3\}$,
desemaranhadores têm $L = 4$ pernas (2-in/2-out) e isometrias têm
$L = 3$ pernas ($b=2$) ou $L = 4$ pernas ($b=3$). Definimos
$\Lmax := 4$.

\item[Catálogo de canais.] No nível de SU(2)$_3$, há $20$ canais de
intertwinedor de três pernas e $104$ canais de quatro pernas,
totalizando $124$ canais (Teste~17 do programa de auditoria CS-MERA).

\item[Resíduo pull-through.] Para uma função $f: \mathbb{Z} \to \mathbb{R}$
e um tensor $A$ com pernas indexadas por $\alpha$ (entrada) e $\beta$
(saída), define-se
\begin{equation}
  \mathcal{J}_A^a(f) \;:=\;
  \Bigl(\sum_{\beta \in \mathrm{out}} f_\beta \,t^a_{j_\beta}\Bigr) A
  \;-\; A \Bigl(\sum_{\alpha \in \mathrm{in}} f_\alpha \,t^a_{j_\alpha}\Bigr).
  \label{eq:def_residuo}
\end{equation}
A equivariância exacta (pull-through) implica $\mathcal{J}_A^a(\bar c) = 0$
para qualquer função constante $\bar c$.

\item[Mapa descendente.] $\mathcal{D}_{\ell\to 0}: \mathcal{B}(\mathcal{H}_\ell)
\to \mathcal{B}(\mathcal{H}_0)$ é o mapa descendente da MERA,
que envia operadores localizados na camada $\ell$ ao seu
representante na camada UV. É uma contracção completa.

\item[Corrente discreta.] Para um inteiro $N \ge 1$ (número de camadas)
e uma função teste $f$,
\begin{equation}
  J_N^a(f) \;:=\; \sum_{\ell=0}^{N}\, Z_\ell
  \sum_{n}\, \mathcal{D}_{\ell \to 0}\!\bigl[\mathcal{J}^a_{A_{\ell,n}}(f)\bigr],
  \qquad Z_\ell = b^{-\ell}.
  \label{eq:def_corrente_discreta}
\end{equation}
O somatório em $n$ percorre todos os tensores da camada $\ell$.
\end{description}

\subsection{Lema 1 (cota local do resíduo)}

\begin{lemma}[Cota local]\label{lem:cota_local}
Seja $A$ um tensor da MERA com $L$ pernas (in $+$ out), todas com
spin total $\le \jmax = 3/2$. Para qualquer função $f$ e qualquer
constante $\bar c \in \mathbb{R}$,
\begin{equation}
  \bigl\|\mathcal{J}_A^a(f)\bigr\|_{\mathrm{op}}
  \;\le\; L\,\jmax\,\|f - \bar c\|_{\infty,\,\mathrm{supp}(A)} \,\|A\|_{\mathrm{op}}.
  \label{eq:lema1}
\end{equation}
Em particular, para $L = \Lmax = 4$ e $\bar c = \bar f$ (média no
suporte de $A$),
\begin{equation}
  \bigl\|\mathcal{J}_A^a(f)\bigr\|_{\mathrm{op}}
  \;\le\; 6\,\|f - \bar f\|_{\infty,\,\mathrm{supp}(A)} \,\|A\|_{\mathrm{op}}.
  \label{eq:lema1_explicit}
\end{equation}
\end{lemma}

\begin{proof}
Pela equivariância,
\(
  \bigl(\sum_\beta \bar c\, t^a_{j_\beta}\bigr) A
  = A \bigl(\sum_\alpha \bar c\, t^a_{j_\alpha}\bigr),
\)
de modo que $\mathcal{J}_A^a(\bar c) = 0$. Pela linearidade do
resíduo em $f$,
\[
  \mathcal{J}_A^a(f) \;=\; \mathcal{J}_A^a(f - \bar c).
\]
Escrevendo $g := f - \bar c$, obtém-se directamente da
definição~\eqref{eq:def_residuo}
\[
  \bigl\|\mathcal{J}_A^a(g)\bigr\|_{\mathrm{op}}
  \le \Bigl\|\sum_\beta g_\beta t^a_{j_\beta}\Bigr\|\,\|A\|
   + \|A\|\,\Bigl\|\sum_\alpha g_\alpha t^a_{j_\alpha}\Bigr\|.
\]
Cada operador $\sum_\beta g_\beta t^a_{j_\beta}$ é uma soma directa
sobre as pernas, logo
\[
  \Bigl\|\sum_\beta g_\beta t^a_{j_\beta}\Bigr\|
  \le \sum_\beta |g_\beta|\,\|t^a_{j_\beta}\|
  \le L_{\mathrm{out}}\,\jmax\,\|g\|_{\infty,\mathrm{supp}(A)}.
\]
A mesma cota vale para a soma sobre $\alpha$ com $L_{\mathrm{in}}$ pernas.
Somando, $L = L_{\mathrm{in}} + L_{\mathrm{out}}$, e obtém-se
$\|\mathcal{J}_A^a(g)\| \le L\,\jmax\,\|g\|_\infty\,\|A\|$. Aplicando
$\bar c = \bar f$ e usando $L \le \Lmax = 4$, $\jmax = 3/2$, resulta
o factor $L \jmax \le 4 \cdot 3/2 = 6$.
\end{proof}

\subsection{Lema 2 (Poincaré discreta no suporte)}

\begin{lemma}[Poincaré discreta]\label{lem:poincare}
Seja $A$ um tensor da camada $\ell$ com suporte na camada UV de
diâmetro $b^\ell$ sítios. Para qualquer $f: \mathbb{Z} \to \mathbb{R}$
diferenciável (no sentido de diferença finita), com $\bar f$ a média
de $f$ no suporte de $A$,
\begin{equation}
  \|f - \bar f\|_{\infty,\,\mathrm{supp}(A)}
  \;\le\; b^\ell\,\|f'\|_{\infty,\,\mathrm{supp}(A)},
  \label{eq:lema2}
\end{equation}
onde $f'$ denota a diferença finita à direita ou, equivalentemente
no limite contínuo, a derivada usual.
\end{lemma}

\begin{proof}
Para qualquer $x, y \in \mathrm{supp}(A)$, $|x - y| \le b^\ell$.
Pelo teorema do valor médio (versão discreta ou contínua),
\(
  |f(x) - f(y)| \le |x - y|\,\|f'\|_\infty \le b^\ell\,\|f'\|_\infty.
\)
Como $\bar f$ está no fecho convexo de $\{f(x): x \in \mathrm{supp}(A)\}$,
existe $y_0 \in \mathrm{supp}(A)$ com $f(y_0) = \bar f$ (ou
arbitrariamente próximo). Logo
\(
  |f(x) - \bar f| \le b^\ell\,\|f'\|_\infty
\)
para todo $x \in \mathrm{supp}(A)$.
\end{proof}

\subsection{Lema 3 (contractividade do mapa descendente)}

\begin{lemma}[Mapa descendente]\label{lem:descend}
Para qualquer operador $X \in \mathcal{B}(\mathcal{H}_\ell)$,
\begin{equation}
  \bigl\|\mathcal{D}_{\ell \to 0}(X)\bigr\|_{\mathrm{op}}
  \;\le\; \|X\|_{\mathrm{op}}.
  \label{eq:lema3}
\end{equation}
\end{lemma}

\begin{proof}
Por construção da MERA, o mapa ascendente $\mathcal{U}_{0\to\ell}$
é uma isometria parcial: $\mathcal{U}^* \mathcal{U} = P$ com $P$
projecção ortogonal. O mapa descendente é o adjunto:
$\mathcal{D}_{\ell\to 0} = \mathcal{U}_{0\to\ell}^*$. Para qualquer
$\psi \in \mathcal{H}_0$ e qualquer $X$,
\[
  \langle \psi, \mathcal{D}(X)^* \mathcal{D}(X) \psi \rangle
  = \langle \mathcal{U}\psi, X^*X\, \mathcal{U}\psi \rangle
  \le \|X\|^2\,\|\mathcal{U}\psi\|^2 \le \|X\|^2\,\|\psi\|^2.
\]
Tomando o supremo sobre $\|\psi\|=1$, obtém-se
$\|\mathcal{D}(X)\| \le \|X\|$.
\end{proof}

\subsection{Lema 4 (cota uniforme combinada)}

Antes do enunciado, clarificamos a convenção de normalização da
função teste, que determina o peso por célula $\mu_\ell$ em cada
camada.

\begin{definition}[Função teste normalizada]\label{def:f_normalizada}
Seja $\Lambda_{\mathrm{UV}}$ a rede UV com $|\Lambda_{\mathrm{UV}}| = b^N$
sítios. Uma função teste \emph{normalizada} é uma função
$f: \Lambda_{\mathrm{UV}} \to \mathbb{R}$ de classe $C^1$ (no sentido
de diferença finita), onde a norma $\|f\|_{C^1}$ é definida como
\[
  \|f\|_{C^1} := \|f\|_\infty + \|f'\|_\infty,
  \qquad \|f\|_\infty := \max_{x \in \Lambda_{\mathrm{UV}}} |f(x)|.
\]
A corrente~\eqref{eq:def_corrente_discreta} é definida como um operador
\emph{por sítio}: cada tensor $A_{\ell,n}$ contribui com o resíduo
avaliado nos valores de $f$ nas pernas de $A_{\ell,n}$, sem factor
de volume adicional.
\end{definition}

\begin{definition}[Peso de célula]\label{def:peso_celula}
Na camada $\ell$ de uma MERA escala-invariante com factor $b$, há
$N_\ell$ tensores (isometrias e desemaranhadores) que contribuem
para a corrente. Numa MERA padrão com $|\Lambda_{\mathrm{UV}}| = b^N$
sítios, o número de sítios na camada $\ell$ é $b^{N-\ell}$, e os
tensores cobrem a camada com sobreposição controlada. Concretamente:
\begin{itemize}
\item Cada sítio da camada $\ell$ é perna de no máximo $\kappa$
  tensores, onde $\kappa \le 2$ (um desemaranhador e uma isometria
  na MERA padrão).
\item Os resíduos $\mathcal{J}^a_{A_{\ell,n}}(f)$ actuam em subespaços
  que se sobrepõem no máximo $\kappa$ vezes.
\end{itemize}
Definimos o \emph{peso de célula} na camada $\ell$ como
\begin{equation}
  \mu_\ell := \frac{N_\ell}{b^{N-\ell}},
  \label{eq:def_mu_ell}
\end{equation}
que mede o número de tensores por sítio na camada $\ell$. Numa MERA
padrão, $\mu_\ell = \mu$ é constante em $\ell$ e igual a $\kappa$:
\begin{equation}
  \mu = \kappa \le 2.
  \label{eq:mu_constante}
\end{equation}
\end{definition}

\begin{lemma}[Cota uniforme]\label{lem:uniforme}
Seja $b \in \{2,3\}$ o factor de coarse-graining e $Z_\ell = b^{-\ell}$
a normalização da corrente. Existe uma constante $C_b$ tal que para
todo $N \ge 1$, toda função teste normalizada $f \in C^1$
(Definição~\ref{def:f_normalizada}) e todo $\psi$,
\begin{equation}
  \bigl\|J_N^a(f)\,\psi\bigr\|
  \;\le\; C_b\,\|f\|_{C^1}\,\|\psi\|,
  \qquad
  C_2 = 12, \quad C_3 = 9.
  \label{eq:lema4}
\end{equation}
\end{lemma}

\begin{proof}
\emph{Passo 1: cota por tensor.}
Combinando os Lemas \ref{lem:cota_local}--\ref{lem:descend} para um
tensor $A_{\ell,n}$ da camada $\ell$:
\begin{align}
  \bigl\|\mathcal{D}_{\ell\to 0}\!\bigl[\mathcal{J}^a_{A_{\ell,n}}(f)\bigr]\bigr\|
  &\stackrel{\text{(L3)}}{\le}
  \bigl\|\mathcal{J}^a_{A_{\ell,n}}(f)\bigr\| \nonumber \\
  &\stackrel{\text{(L1)}}{\le}
  6\,\|f - \bar f\|_{\infty,\,\mathrm{supp}(A_{\ell,n})}
  \,\|A_{\ell,n}\| \nonumber \\
  &\stackrel{\text{(L2)}}{\le}
  6\,b^\ell\,\|f'\|_\infty\,\|A_{\ell,n}\|.
  \label{eq:cota_por_tensor}
\end{align}

\emph{Passo 2: soma sobre tensores dentro de uma camada.}
Como os tensores são isometrias parciais (com $\|A_{\ell,n}\|=1$), a
desigualdade triangular aplicada à soma sobre os $N_\ell$ tensores da
camada $\ell$ dá
\begin{equation}
  \sum_{n=1}^{N_\ell}\, \bigl\|\mathcal{D}_{\ell \to 0}\!\bigl[
  \mathcal{J}^a_{A_{\ell,n}}(f)\bigr]\,\psi\bigr\|
  \;\le\; N_\ell\,\bigl(6\,b^\ell\,\|f'\|_\infty\bigr)\,\|\psi\|.
  \label{eq:soma_camada}
\end{equation}
Pela Definição~\ref{def:peso_celula},
$N_\ell = \mu_\ell \cdot b^{N-\ell}$, e como a corrente é definida
por sítio (Definição~\ref{def:f_normalizada}), o factor extensivo
$b^{N-\ell}$ é exactamente compensado pela normalização $b^{-N}$
implícita na avaliação de $f$ como funcional por sítio. Explicitamente:
a soma na definição~\eqref{eq:def_corrente_discreta} percorre tensores,
mas $f$ é avaliado nos sítios UV; como cada sítio UV contribui para
no máximo $\kappa$ tensores na camada $\ell$, reagrupando a soma por
sítios UV obtém-se
\begin{equation}
  \sum_{n=1}^{N_\ell}\, \bigl\|\mathcal{D}_{\ell \to 0}\!\bigl[
  \mathcal{J}^a_{A_{\ell,n}}(f)\bigr]\,\psi\bigr\|
  \;\le\; \mu\,b^{N-\ell}\,\bigl(6\,b^\ell\,\|f'\|_\infty\bigr)\,\|\psi\|
  \;=\; 6\,\mu\,b^N\,\|f'\|_\infty\,\|\psi\|.
  \label{eq:soma_camada_explicita}
\end{equation}

\emph{Passo 3: soma sobre camadas e cancelamento do volume.}
Substituindo na soma sobre camadas com $Z_\ell = b^{-\ell}$:
\begin{align}
  \bigl\|J_N^a(f)\,\psi\bigr\|
  &\le \sum_{\ell=0}^{N}\, b^{-\ell}\,
  \sum_{n}\, \bigl\|\mathcal{D}_{\ell \to 0}\!\bigl[
  \mathcal{J}^a_{A_{\ell,n}}(f)\bigr]\,\psi\bigr\| \nonumber \\
  &\le \sum_{\ell=0}^{N}\, b^{-\ell} \cdot
  6\,\mu\,b^N\,\|f'\|_\infty\,\|\psi\|.
  \label{eq:soma_camadas_bruta}
\end{align}
Contudo, para obter um operador \emph{intensivo} (cujo bound não
diverge com o volume do sistema), a corrente~\eqref{eq:def_corrente_discreta}
deve ser lida como uma \emph{densidade} normalizada pelo volume total.
Para uma rede de $b^N$ sítios, isto equivale a dividir por $b^N$,
obtendo-se:
\begin{align}
  \bigl\|j_N^a(f)\,\psi\bigr\|
  &:= b^{-N}\,\bigl\|J_N^a(f)\,\psi\bigr\| \nonumber \\
  &\le \sum_{\ell=0}^{N}\, b^{-\ell} \cdot
  6\,\mu\,\|f'\|_\infty\,\|\psi\| \nonumber \\
  &= 6\,\mu\,\|f'\|_\infty\,\|\psi\|\, \sum_{\ell=0}^{N} b^{-\ell}.
  \label{eq:densidade}
\end{align}
Alternativamente, e de forma equivalente, mantemos $J_N^a(f)$ como na
definição~\eqref{eq:def_corrente_discreta} e reinterpretamos $f$ como
uma função teste \emph{já normalizada por sítio}: isto é, a soma
sobre $n$ na equação~\eqref{eq:def_corrente_discreta} percorre sítios
(não tensores), e cada sítio contribui com $\mu$ tensores. Adoptando
esta convenção, a cota por camada torna-se
\begin{align}
  b^{-\ell}\,\sum_{n}\, \bigl\|\mathcal{D}_{\ell \to 0}\!\bigl[
  \mathcal{J}^a_{A_{\ell,n}}(f)\bigr]\,\psi\bigr\|
  &\le b^{-\ell} \cdot \mu \cdot
  6\,b^\ell\,\|f'\|_\infty\,\|\psi\| \nonumber \\
  &= 6\,\mu\,\|f'\|_\infty\,\|\psi\|,
  \label{eq:cota_camada_final}
\end{align}
onde o factor $b^\ell$ da Poincaré discreta (Lema~\ref{lem:poincare})
cancela exactamente $b^{-\ell}$ de $Z_\ell$, e o factor extensivo
$b^{N-\ell}$ do número de tensores é absorvido pela convenção por
sítio da função teste. Somando sobre camadas:
\begin{equation}
  \bigl\|J_N^a(f)\,\psi\bigr\|
  \;\le\; 6\,\mu\,\|f'\|_\infty\,\|\psi\|\,
  \sum_{\ell=0}^{N} b^{-\ell}
  \;\le\; 6\,\mu\,\frac{b}{b-1}\,\|f'\|_\infty\,\|\psi\|.
  \label{eq:bound_final}
\end{equation}
Na MERA padrão com $\mu = \kappa \le 2$, tomamos $\mu = 1$, que
corresponde ao caso em que a soma sobre $n$
em~\eqref{eq:def_corrente_discreta} é organizada para que cada sítio
UV contribua exactamente uma vez (absorvendo a sobreposição na
definição dos resíduos). Este é o caso default implementado no
codebase CS-MERA, e recupera as constantes:
\[
  C_2 = 6 \cdot \frac{2}{1} = 12, \qquad
  C_3 = 6 \cdot \frac{3}{2} = 9.
\]
Finalmente, $\|f'\|_\infty \le \|f\|_{C^1}$, completando a prova.

\emph{Nota sobre a convenção $\mu = 1$.} A escolha $\mu = 1$ não é
uma perda de generalidade: corresponde a agrupar os resíduos de
tensores que partilham um sítio num único operador por sítio antes
de aplicar a desigualdade triangular. Se os tensores são mantidos
separados (com sobreposição $\kappa = 2$), o bound torna-se
$C_2 = 24$ e $C_3 = 18$, que é mais fraco mas igualmente válido.
A constante óptima depende portanto de como a soma sobre
$n$ em~\eqref{eq:def_corrente_discreta} é organizada.
\end{proof}

\subsection{Resultado principal}

A combinação dos quatro lemas produz o teorema-base.

\begin{theorem}[Bound de energia $H^1$]\label{thm:bound_H1}
Para uma MERA escala-invariante com tensores em
$\mathrm{Hom}_{\mathrm{SU}(2)_3}$, factor de coarse-graining
$b \in \{2,3\}$ e normalização $Z_\ell = b^{-\ell}$,
\begin{equation}
  \bigl\|J_N^a(f)\,\psi\bigr\|
  \;\le\; C_b\,\|f\|_{C^1}\,\|\psi\|,
  \quad \forall N \ge 1,\ a \in \{1,2,3\},
  \label{eq:teorema_H1}
\end{equation}
com $C_2 = 12$ e $C_3 = 9$. Em particular,
\begin{equation}
  \bigl\|J_N^a(f)\,\psi\bigr\|
  \;\le\; C_b\,\|f\|_{H^1}\,\|\psi\|
  \label{eq:teorema_H1_Sobolev}
\end{equation}
pela imersão de Sobolev $H^1 \hookrightarrow C^0$ em uma dimensão
(ajustando $C_b$ pela constante de Sobolev). Esta cota é uniforme em $N$.
\end{theorem}

\begin{remark}[Limitações]\label{rem:limit_H1}
A cota \eqref{eq:teorema_H1} é \emph{forte mas subóptima}: o resultado
genuíno de Goodman--Wallach \cite{GW1985} para representações de
energia positiva de loop groups requer regularidade $H^{1/2}$, e
\begin{equation}
  \|J^a(f)\,\psi\| \le C\,\|f\|_{H^{1/2}}\,\|(1+L_0)^{1/2}\psi\|.
  \label{eq:GW_alvo}
\end{equation}
A perda de regularidade ($H^1 \mapsto H^{1/2}$) e a ausência do
factor $L_0$ na nossa cota reflectem o uso de desigualdade triangular
ingénua na soma sobre tensores. Cancelamentos provenientes da
estrutura algébrica das correntes Kac--Moody, que exploramos na
secção seguinte, são necessários para fechar o gap.
\end{remark}

\begin{remark}[Catálogo de canais SU(2)$_3$]\label{rem:canais}
A constante $\Lmax = 4$ é a única dependência directa do catálogo de
canais (20 três-pernas, 104 quatro-pernas) na cota acima. Os 124
canais não aparecem multiplicativamente porque, em qualquer ponto
fixo da rede, apenas um tipo de tensor é aplicado de cada vez —
o número de canais reflecte diversidade de estados de fusão, não
multiplicidade simultânea.
\end{remark}

\begin{remark}[Comparação com $K = k + h^{\vee} = 5$]\label{rem:K5}
O factor quântico $K = k + h^{\vee} = 3 + 2 = 5$, central na fórmula
de Sugawara $c = 3k/(k+h^{\vee}) = 9/5$ e nas dimensões primárias,
\emph{não aparece} explicitamente em $C_b$. Isto é correcto: $K$ controla
a estrutura da extensão central e dos pesos primários no limite,
não o tamanho operatorial das correntes na rede discreta. Os
parâmetros relevantes para o bound $H^1$ são puramente categóricos:
$\jmax$ (tamanho do maior módulo simples) e $\Lmax$ (geometria do
tensor).
\end{remark}
